% Lunettes à anaglyphes : ce que voit chaque œil (corrigé de l'exercice 13, chapitre 19)
\definecolor{pcR}{RGB}{230,30,30}
\begin{tikzpicture}[x=1cm,y=1cm]
% Image imprimée
\begin{scope}
  \draw[otRule,line width=0.5pt] (0,0) rectangle (2.4,1.6);
  \draw[pcR,line width=2pt] (0.5,0.4) rectangle (1.3,1.1);
  \draw[cyan!80!black,line width=2pt] (0.8,0.5) rectangle (1.6,1.2);
  \node[pcLegendeG,anchor=north] at (1.2,-0.08) {image imprimée};
\end{scope}
% Oeil gauche : filtre rouge
\begin{scope}[shift={(3.4,0)}]
  \fill[pcR!85] (0,0) rectangle (2.4,1.6);
  \draw[black,line width=2pt] (0.8,0.5) rectangle (1.6,1.2);
  \node[pcLegendeG,anchor=north,align=center] at (1.2,-0.08) {œil gauche (filtre rouge) :\\le dessin cyan, noir};
\end{scope}
% Oeil droit : filtre cyan
\begin{scope}[shift={(6.8,0)}]
  \fill[cyan!80] (0,0) rectangle (2.4,1.6);
  \draw[black,line width=2pt] (0.5,0.4) rectangle (1.3,1.1);
  \node[pcLegendeG,anchor=north,align=center] at (1.2,-0.08) {œil droit (filtre cyan) :\\le dessin rouge, noir};
\end{scope}
\end{tikzpicture}
\par\smallskip\legende{À travers chaque filtre, le fond blanc et le dessin de la couleur du filtre se confondent ; seul l'autre dessin apparaît, en noir.}
